% Generated draft: requires mathematical and source review.
\documentclass[11pt]{article}
\usepackage[margin=25mm]{geometry}
\usepackage{fontspec}
\setmainfont{Latin Modern Roman}
\usepackage{amsmath,amssymb,mathtools,unicode-math}
\setmathfont{Latin Modern Math}
\usepackage{enumitem,needspace,xurl,hyperref}
\setlength{\parindent}{0pt}
\setlength{\parskip}{5pt}
\setlength{\emergencystretch}{3em}
\newcommand{\N}{\mathbb{N}}
\newcommand{\Z}{\mathbb{Z}}
\newcommand{\Q}{\mathbb{Q}}
\newcommand{\R}{\mathbb{R}}
\newcommand{\C}{\mathbb{C}}
\newcommand{\F}{\mathbb{F}}
\newcommand{\A}{\mathbb{A}}
\newcommand{\PP}{\mathbb{P}}
\newcommand{\E}{\mathbb{E}}
\newcommand{\eps}{\varepsilon}
\newcommand{\dd}{\,\mathrm{d}}
\begin{document}
\section*{UnsolvedMath ID: 20001004}
\textbf{Source problem code:} catalogue entry
\subsection{Definitions and mathematical statement}
Let $Y$ be a finite connected loopless multigraph with Betti genus $g=|E(Y)|-|V(Y)|+1\ge2$. The original question, posed by Gregg Musiker, asks to recognize when this given graph admits a harmonic automorphism group of order $6(g-1)$ in Corry's sense. Scott Corry's accompanying remark concerns constructing a graph from a known maximal group by a Cayley-type method; Musiker explicitly distinguishes that construction from recognition of an input graph.

A dart is an edge together with a choice of initial endpoint, so $|D(Y)|=2|E(Y)|$. The attached aid states that recognition is equivalent to existence of $G\le\operatorname{Aut}(Y)$ of order $6(g-1)$ acting semiregularly on darts, meaning every dart stabilizer is trivial. For multigraphs the automorphism action includes the actual edge incidences, not merely vertex permutations. This supplies a finite subgroup-search test; no polynomial complexity bound is claimed. The exact harmonic-quotient convention from Corry [15] is not reproduced and must be aligned with this criterion.

A necessary arithmetic consequence, including for graphs with leaves, is
\[
|E(Y)|=3k(g-1)\quad\text{for an integer }k\ge1.
\]
The aid further attributes a stronger leafless reduction to Corry's uniform-subdivision theorem. If the minimum degree is at least two, it states that maximality is equivalent to $Y$ being a uniform subdivision of a loopless cubic graph $C$, all core edges replaced by paths of the same positive length, with the following rotation certificate.

Choose a cyclic order of the three darts at every vertex of $C$. Let $\rho$ be the product of these $3$-cycles and let $\tau$ reverse each dart across its edge. The condition is
\[
|\langle\rho,\tau\rangle|=|D(C)|.
\]
Connectedness makes this monodromy group transitive, so the equality says it is regular. Its centralizer in $\operatorname{Sym}(D(C))$ is then also dart regular and gives the corresponding graph automorphism action. The aid states existence of some such rotation choice exactly characterizes admissible cores. Suppressing degree-two vertices must in particular produce a loopless cubic core and uniform path lengths; an arbitrary suppression need not satisfy either property.

This is the aid's precise fixed-graph certificate and reduction, rather than a newly verified complete graph classification. Graphs with leaves retain the general subgroup test and divisibility obstruction; the leafless rotation certificate is not asserted for them.
\subsection{Short English statement}
Recognize maximal harmonic symmetry on a specified graph. The supplied aid reduces the general case to a free dart action of the required order and the leafless case to a uniform cubic subdivision with a finite rotation-group certificate.
\subsection{Sources}
\begin{itemize}
\item Local UnsolvedMath registry, permanent ID $20001004$, source code catalogue entry; canonical \texttt{aim-combinatorics-notes.json}, \texttt{notes[128]}.
\item American Institute of Mathematics, 2013 ``Generalizations of chip-firing and the critical group,'' problem 39, posed by Gregg Musiker with a response by Scott Corry: \url{https://aimath.org/pastworkshops/chipfiringproblems.pdf}. The source cites Corry as [15], without its full bibliographic entry here. No PDF inspection is claimed.
\item Attached research attempt 1: partial result, run and registry status partially solved, rigor proved, novelty confidence low, classification PARTIAL-PROGRESS. It invokes Corry's uniform-subdivision theorem without an exact citation. The machine-generated aid is not peer reviewed; independently verify before citation.
\end{itemize}
\end{document}
